% Section content template for individual Educative lessons
% This file contains the actual content for a specific section
% Generated from Educative API section content

% Section: Quiz on the Distributed Messaging Queue’s Design
% Section ID: 5261732032806912
% Section Slug: quiz-on-the-distributed-messaging-queues-design
% Generated: 2025-12-07T20:58:30.453101

\subsection*{Quiz}
\vspace{0.3cm}
\noindent
\textbf{Question 1:} What kind of messages do dead-letter queues contain?
\\[0.2cm]
\begin{itemize}
 \item The messages that have been consumed successfully.
 \item \textbf{[CORRECT]} The messages that failed to be consumed and have reached the maximum attempts limit.
\\[0.1cm]
\textit{Explanation:} 
A dead letter cannot be delivered or returned because it lacks appropriate directions. Similarly , a queue named after the \texttt{dead-letter} contains those messages that have failed to be consumed after the maximum attempts limit.
 \item The messages that have been produced by a process that's dead now.
 \item None of the above
\end{itemize}
\vspace{0.3cm}
\noindent
\textbf{Question 2:} The purpose of replicating queues on multiple servers is to enhance the \_\_\_\_\_\_\_\_\_ of the system.
\\[0.2cm]
\begin{itemize}
 \item security
 \item consistency
 \item \textbf{[CORRECT]} availability
\\[0.1cm]
\textit{Explanation:} 
The purpose of replicating the queues and corresponding messages on multiple nodes is to increase the system's availability. When one node is down, another stand-by server can serve the request.
 \item None of the above
\end{itemize}
\vspace{0.3cm}
\noindent
\textbf{Question 3:} In the \textbf{cluster of independent hosts} model, which component is responsible for replicating messages in the other nodes?
\\[0.2cm]
\begin{itemize}
 \item The front-end server
 \item The cluster manager
 \item \textbf{[CORRECT]} Any random host within the cluster
\\[0.1cm]
\textit{Explanation:} 
In the \textbf{cluster of independent hosts} model, a random host can receive the message responsible for replicating messages in other hosts in the respective queues.
 \item None of the above
\end{itemize}
\vspace{0.3cm}
\noindent
\textbf{Question 4:} The purpose of a metadata service is \_\_\_\_\_\_\_\_\_.
\\[0.2cm]
\begin{itemize}
 \item to store messages in the queues along with their metadata
 \item \textbf{[CORRECT]} to store, retrieve, and update the metadata of queues in the metadata store and cache
\\[0.1cm]
\textit{Explanation:} 
Metadata service acts as a middleware between the front-end and metadata storages (and caches). It is responsible for storing, retrieving, and updating the metadata of queues in the metadata stores and caches.
 \item to provide security features in accessing the queues
 \item None of the above
\end{itemize}
\vspace{0.3cm}
\noindent
\textbf{Question 5:} Which of the following is responsible for partitioning the queue and assigning a primary node to each partition?
\\[0.2cm]
\begin{itemize}
 \item \textbf{[CORRECT]} The internal cluster manager
\\[0.1cm]
\textit{Explanation:} 
The internal cluster manager partitions a queue into multiple parts and assigns a primary node to each partition. In contrast, the external cluster manager assigns clusters to each queue partition.
 \item The external cluster manager
 \item The primary node
 \item One of the secondary nodes
\end{itemize}

% End of section content